\cotefile{79}
\scanpage{81}
les structures $(\widetilde F_{M,i}\ \unsure{\text{de}\mathbin{+}\pi|\text{déf}\ \pi|)_{i\in\pi}|)_{i\in\Pi}}$.

Soit $u$ un autom. de cette structure,
commutant aux \deleted{\unsure{K}} $\widetilde T_\lambda$, et fixant un
repère $(x,y,z)$, i.e. fixant les éléments
d'une base $(x,y)$.
% slanted
\margin{il suffit pour voir qu'il fixe $x,y$} \deleted{Alors} Notons que un
tel élément $\xi$, \inserted{$\in\widetilde S_M=M^{**}$} il fixe les $T_\lambda\xi=\lambda\xi$, $\lambda\in\Lambda^*$,
en particulier aussi $-\xi$ -- \deleted{on voit donc}
\inserted{Comme il} \deleted{par le lemme p.~$37'$, que \unsure{pour} $a,b\in$}
\inserted{si $x$} fixe $x$ et les $\lambda y$ $(\lambda\in\Lambda^*)$, on voit que il
\deleted{induit l'identit\unsure{é}} fixe les $\tau_A(\{x,\lambda y\})=\lambda\cdot i$
$(\text{où}\ i=\deleted{\unsure{\partial}}\tau_A\{x,y\})$, donc
il induit l'identité sur $\pi_M$, i.e. \deleted{\unsure{Raisonnons}} \inserted{c'est un}
\deleted{\unsure{P}}autom. pour toutes les structures $\widetilde F_{M,i}$.

\deleted{En vertu du lemme p.~$37'$, avec $x,y$,}
\deleted{Il fixe aussi} Considérons les matrices
\deleted{$u$ fixe aussi les $ax+by$ avec $a$}
\[
\left\{\begin{aligned}
\underline{\lambda}_0(\lambda)&=\begin{pmatrix}1&-\lambda\\0&1\end{pmatrix}
&&\lambda\in\Lambda^*\\
\lambda_1&=\begin{pmatrix}0&1\\-1&0\end{pmatrix}\\
\lambda_\infty&=\begin{pmatrix}1&-1\\1&0\end{pmatrix}
\end{aligned}\right.
\]

Alors l'ens. des bases $(x,y)$ ``fixes sous $u$'' est stable par composition avec les matrices
précédentes et leurs inverses, donc aussi
les composés de telles. Si $\Lambda^*$ engendre $\Lambda$
additivement (p.~ex. $\Lambda$ $\frac12$ local) alors \deleted{\unsure{du}} être
% archive sheet number: 80
\scanpage{82}\mspage{40}
peut prendre $=\lambda_0(\lambda)$ avec $\lambda\in\Lambda$ quelconque (pas néc$^{\text{t}}$ $\lambda\in\Lambda^*$). Mais sauf erreur,
ces trois types de matrices engendrent
$\mathrm{Sl}(2,\Lambda)$. On trouve ainsi (sauf erreur!)
% slanted
\margin{à vérifier}
que $u$ est l'identité sur $\widetilde S_M$ \deleted{tout} entier,
\deleted{vu que $\mathrm{Sl}(2,\Lambda)$ est transitif}
Comme $\mathrm{Gl}_\Lambda(M)$ est transitif sur l'ens. des
repères, on trouve ainsi :

Si $\Lambda^*$ engendre $\Lambda$ additivement, \deleted{alors} p.~ex. $\Lambda$ anneau
$\frac12$ local, alors tout autom. de la structure
\deleted{$\unsure{\widetilde P_M|\widetilde F_M}$} $(\widetilde S_M,\underset{\unsure{\mathfrak P(\mathfrak P_3(S))|\mathbb P(\mathbb P_3(S))|\mathcal P(\mathcal P_3(S))}}{\underset{\rotatebox[origin=c]{-90}{\(\in\)}}{\widetilde F_M}},T_\bullet:\Lambda^*\longrightarrow\circlearrowright_{\widetilde S_M})$ provient d'un
autom. unique de $M$.

Notons que pour une permutation
$u$ de $\widetilde S_M=M^*$, dire qu'elle commute
à l'antipodisme et commute à la structure
\unsure{triangulée|triangulation|triangulaire} signifie que l'on a
\[
\left\{\begin{aligned}
\text{a)}\ u(-x)&=-u(x)&&x\in M^*\\
\text{b)}\ u(x+y)&=u(x)+u(y)&&
\begin{aligned}&x,y\,\deleted{\unsure{x+y}}\in M^*\\[-2pt]&x\wedge y\in(\det M)^*\end{aligned}
\end{aligned}\right. \tag*{(21)}
\]
et dire qu'elle commute de plus à
l'action de $\Lambda^*$ signifie qu'on a, au
lieu de $(21\text{ a})$, la condition plus forte
\[
u(\lambda x)=\lambda u(x)\qquad x\in M^*,\ \lambda\in\Lambda^*. \tag*{(22)}
\]
% archive sheet number: 81
\scanpage{83}
Pour montrer que $u$ est induit par un
autom. (unique) de $M$, il est tentant
d'écrire $M$, en termes de \deleted{l'action}
\deleted{de} $(\widetilde S_M,\widetilde F_M,T_\bullet:\Lambda^*\,-\,\unsure{\circlearrowright|\mathfrak S}_{\widetilde S_M})$, comme
\[
\begin{aligned}
M\xleftarrow{\sim}{}&\text{quotient du \inserted{$\mathbb Z$-}module libre }\mathbb Z^{(M^*)}\\
&\text{par le sous-groupe engendré}\\
&\text{par les }\ [x+y]-[x]-[y]\\
&\text{(pour }x,y,\deleted{\unsure{x+y}}\in M^*\ \text{tels que }x\wedge y\in(\det M)^*)\\
&\text{et par les }\ [-x]+[x]\;.
\end{aligned}\tag*{(23)}
\]
% slanted
\margin{\uline{NB} La structure ``triangulation'' $+$ ``antipod.'' peut s'interpréter (en termes de l'axiome de l'unicité du triangle complétant une arête) par une loi de groupe comm. pas partout définie $(x,y)\mapsto x\unsure{\mathbin{\perp}|\mathbin{+}}y$, $x\mapsto -x$, dont $M$ pourrait alors s'interpréter comme la ``\unsure{quasi-addition|quasi-loi|groupe abélien|groupe}'' \unsure{enveloppante|enveloppant}\ldots}

Désignant par $R_M$ ce sous-groupe,
une suite d'homomorphismes can.
\diagram{\(\text{(24)}\quad\begin{tikzcd}[ampersand replacement=\&,column sep=1.3em]0\arrow[r]\&R_M\arrow[r,"i"]\&\mathbb Z^{M^*}\arrow[r,"p"]\&M\arrow[r]\&0\end{tikzcd}\)}
et il faudrait montrer qu'elle est
exacte (le seul pb \unsure{ici|il} est $\operatorname{Ker}p=\operatorname{Im}i$,
alors qu'il est clair que $pi=0$). (Une
question analogue peut se poser pour un
module libre de rang quelconque \inserted{$\geq 2$}.)

Ainsi, $u$ satisfaisant (21) induit une
\deleted{application} additif $\widehat u$ de $M$ qui l'induit,
et la condition (22) implique que ce
dernier commute à $\Lambda^*$, donc à $\Lambda$ si $\Lambda^*$
engendre $\Lambda$ additivement\ldots{} Il est clair a
% archive sheet number: 82
\scanpage{84}\mspage{41}
priori que cette condition supplémentaire
est inutile \deleted{\unsure{q}}pour impliquer que $\widehat u$ est
$\Lambda$-linéaire, quand $\Lambda$ est un anneau
absolu $(\simeq\mathbb Z\ \text{ou}\ \mathbb Z/n\mathbb Z)$, puisque tt
autom. additif d'un $\Lambda$-module est
alors au\deleted{\unsure{toma|tomat}}tomatiquement $\Lambda$-linéaire. Par
contre on comprend bien pourquoi,
quand $\Lambda$ n'est pas un anneau absolu,
il faut absolument mettre une structure
supplémentaire sur $\widetilde P_M$ \inserted{$=(\unsure{S_M|\widetilde S_M},\widetilde F_M)$} (genre op. de $\Lambda^*$
dessus) -- en effet tout autom. \uline{additif}
(pas néc. $\Lambda$-linéaire) de $M$ induit un
autom. de $\widetilde P_M$ ! Mais cette remarque
ne s'applique pas telle quelle, si on
tient compte de la structure suppl.
$\widetilde F_M=\coprod_i\widetilde F_{M,i}$ -- \deleted{sauf si} \inserted{si pourtant \unsure{si|si si}} $\Lambda^*=\{\pm1\}$
(p.~ex. $\Lambda=$ \deleted{$\mathbb F_3$} anneau de polynômes sur
$\mathbb F_2,\mathbb F_3$ ou $\mathbb Z$), auquel cas l'ens. d'indices
$\pi$ lui-m$^{\text{me}}$ est réduit à un élément !
% archive sheet number: 83
\scanpage{85}
J'en arrive maintenant au cas \uline{B})
et je supposerai dorénavant, pour ce cas
\[
\boxed{2\cdot1_\Lambda\ne0}\tag*{(25)}
\]
de façon que \inserted{l'involution} $-\mathrm{id}_M$ opère librement sur
$\widetilde S_M=M^*$, sur $\widetilde A_M$ et sur $\widetilde F_M$ -- on
note \deleted{que} effet que
\[
\begin{aligned}
\{x,y,z\}\in\widetilde F_M\ \Longrightarrow\ &x,y,z,-x,-y,-z\ \text{sont}\\
&\text{tous distincts}.
\end{aligned}\tag*{(26)}
\]

On pose
\[
\left\{\begin{aligned}
S_M&=\widetilde S_M/\{\pm\mathrm{id}_M\}=M^*/\{\pm1\}\\
\deleted{A}F_M&=\deleted{\widetilde A}\widetilde F_M/\{\pm\mathrm{id}_M\}=
\end{aligned}\right.\tag*{(27)}
\]
et on a une bijection canonique
\[
F_M\xrightarrow{\sim}\left\{f\in\unsure{\mathfrak P|\mathbb P|\mathcal P}_3(S_M)\mid
\begin{aligned}&\exists\,\widetilde f\in\widetilde F_M,\ \text{avec}\\
&f=p(\widetilde f)\end{aligned}\right\}.\tag*{(28)}
\]
où
\[
\begin{aligned}p:\widetilde S_M&\longrightarrow S_M&\quad&\text{est l'application}\\
x&\longmapsto\dot x&&\text{canonique}.\end{aligned}\tag*{(29)}
\]

On vérifie \deleted{en effet} que si $f\in\mathfrak P_3(S_M)$
satisfait (28), alors il y a exactement deux
$\widetilde f$ tels que $f=p(\widetilde f)$, formant une orbite sous
$\{\pm\mathrm{id}_M\}$.
% archive sheet number: 84
\scanpage{86}\mspage{42}
Par la suite on identifie tjrs $F_M$ à une
partie de $\unsure{\mathfrak P|\mathbb P|\mathcal P}_3(S_M)$.

On \deleted{\unsure{pose}} trouve, par passage aux quotients
\deleted{(30)} à partir \deleted{\unsure{A}} de $\overrightarrow{\widetilde\tau}_F$ \deleted{$\widetilde F_M$} et $\widetilde\tau_F$ (5) (6)
des applications canoniques
\[
\begin{aligned}
\underset{\substack{\text{faces}\\\text{orientées}}}{\overrightarrow F_M}
&\xrightarrow{\overrightarrow\tau_F}\overrightarrow\pi_M=(\det M)^*\\
&&\overrightarrow\tau_F(\check f)=-\overrightarrow\tau_F(f)
\end{aligned}\tag*{(30)}
\]
\[
F_M\xrightarrow{\tau_F}\pi_M=(\det M)^*/\{\pm1_\Lambda\},\tag*{(31)}
\]
ce qui, pour $i\in\pi_M$, permet de définir
les $F_{M,i}$, d'où
\[
F_M=\coprod_{i\in\pi_M}F_{M,i}.\tag*{(32)}
\]

On pose
\[
A_M=\{a\in\mathfrak P_2(S_M)\mid\exists f\in S_M,\ f\supset a\}\tag*{(33)}
\]
et on trouve que l'on a aussi
\[
A_M=\left\{a\in\mathfrak P_2(S_M)\mid
\begin{aligned}&\exists\widetilde a\in\widetilde A_M,\ \text{avec}\\
&a=p(\widetilde a)\end{aligned}\right\}\tag*{(34)}
\]
-- mais ici le nb des $\widetilde a$ \inserted{$\in\widetilde A_M$} au dessus de $a$
est égal à 4 -- si $a=\{s,t\}$, et
$s=\dot x,t=\dot y$, les quatre choix sont
$\{x,y\},\{-x,-y\}\deleted{\unsure{\{x,y\}_M}},\{x,-y\},\{-x,y\}$
% archive sheet number: 85
\scanpage{87}
qui se groupent en deux paires \deleted{d'arêtes}
\deleted{\unsure{de deux}} formées d'arêtes antipodiques de $\widetilde{P}_M$.
Ces deux paires correspondent aux
deux faces de \deleted{$\unsure{F|P}$}
\[
P_M=(S_M,F_M)
\]
qui contiennent $a$ :
\[
\begin{aligned}
&\text{Pour toute }a\in A_M,\ \text{il existe exactement}\\
&\text{\uline{deux} faces } f\in F_M\ \text{telles que } f\supset a.
\end{aligned}\tag*{(35)}
\]

On a de plus, pour mémoire :
\[
\left\{\begin{aligned}
&\text{Soit } e_{\deleted{\unsure{M}}}\in\overrightarrow\pi,\quad i=\{\pm e_{\deleted{M}}\},\quad\text{soit }\deleted{\unsure{F_+}}\\
&\overrightarrow F_{M,e_{\deleted{M}}}^{\deleted{\unsure{+}}}\subset\overrightarrow F_{M,i}\\
&\shortparallel\\[-3pt]
&\{\vec f\in\overrightarrow F_{M,i}\mid\vec\tau(\vec f)=e_{\deleted{M}}\}\\
&\text{alors }F_{M,e}\text{ est une ``orientation'' de }P_{M,i},\\
&\text{i.e. si }\vec f,\vec f'\in\overrightarrow F_{M,e}\text{ sont \uline{adjacentes}}\\
&\text{i.e. }|\vec f|\cap|\vec f'|=a\in\unsure{\mathfrak P|\mathbb P|\mathcal P}_2(S_M),\ \text{alors les}\\
&\text{orientations de }a\text{ induites par }\vec f,\vec f'\text{ sont}\\
&\text{opposées}.
\end{aligned}\right.\tag*{(36)}
\]
% slanted
\margin{\uline{NB} Si on interprète $i$ comme un ens. à 2. éléments \unsure{$(\subset\overrightarrow\pi)$|$(\in\overrightarrow\pi)$|$(\subset\overrightarrow\Pi)$}, on a $\underset{\inserted{\text{can}}}{\operatorname{Or}}(P_{M,i})\simeq i$}

On récupère les $F_{M,i}$ à partir des $A_{M,i}$
(mais non $F_M$ à partir de $A_M$) par
\[
F_{M,i}=\{f\in\mathfrak P_3(S_M)\mid\forall a\in\mathfrak P_2(f),\ \text{on a }a\in A_{M,i}\}.\tag*{(37)}
\]
% archive sheet number: 86
\scanpage{88}\mspage{43}
\deleted{Pour étudier l'étoile d'un}
La propriété (35) montre que les \deleted{$P_{M,i}=$}
\[
P_{M,i}=(S_M,A_{M,i},F_{M,i}\inserted{+\text{ relations d'incidence (\unsure{inclusion|induites|l'inclusion})}})\qquad(i\in\pi)\tag*{(38)}
\]
sont des structures de 2-polyèdres triangulés, sauf qu'il n'est pas clair que l'étoile de chaque sommet soit connexe
(en tant que 1-contour). Choisissons
un $e\in i=\{\pm e\}\subset\overrightarrow\pi=\det(M)^*$, on
trouve que le contour combinatoire orienté
correspondant à un sommet $s=\dot x$ est
donné par
\[
\begin{aligned}
\underbrace{\overrightarrow{\mathrm{Et}}_e(s)}_{\substack{\text{étoile}\\\text{orientée,}\\\text{définie par l'}\\\text{orientation }e\\\text{de }P_{M,i}}}
\simeq{}&\text{contour défini par l'espace}\\
&\text{affine }L_x\subset M\ \text{formé des}\\
&y\text{ tels que }x\wedge y=e,\\
&\text{ayant comme module des}\\
&\text{translations }V_x=\{u\in M\mid x\wedge u=0\}\\
&\hspace{6em}=\Lambda x,\\
&\text{avec comme ens. de}\\
&\text{sommets }L_x,\ \text{et le successeur}\\
&\text{de }y\in L_x\text{ étant }y+\deleted{\unsure{x|1}}\unsure{1|x}
\end{aligned}\tag*{(39)}
\]
% slanted
\margin{\uline{NB} \deleted{\unsure{P|Pour}} A tout $s\in S_M$, le choix de $e\in i$ définit sur \deleted{est associé \unsure{un $\Lambda$|son $\Lambda$}} $\mathrm{\acute Et}_e(s)$ une structure de \uline{torseur} \deleted{module libre de rang 1} \uline{sous $\Lambda$} (de telle façon que le \deleted{$V_s$ de $M$, $V_s=\Lambda\cdot\{x\}=\Lambda\cdot(-x)$} successeur de $y$ soit $y+1$), \deleted{(où $s=\dot x$), et $\mathrm{\acute Et}_e(s)$ est} qui est remplacée par la structure opposée si $e$ remplacé par $-e$. Donc $\mathrm{\acute Et}(s)$ est un torseur sous $\Lambda\wedge_{\mu_2}i$ ($\Lambda$ tordu par $i$)}

On trouve donc
\[
\begin{aligned}
&\text{Les étoiles des points de }P_{M,i}\text{ sont connexes}\\
&\deleted{\unsure{\text{ssi}}}\text{ ssi on a }\Lambda_0=\Lambda\quad\text{i.e. }\Lambda\text{ est un anneau absolu}\ ;\\
&\Lambda\simeq\mathbb Z\text{ ou }\Lambda\simeq\mathbb Z/n\mathbb Z.\quad\text{Dans \deleted{ce cas} \inserted{tous les cas}, les}\\
&\text{ordres des \deleted{\unsure{points}}sommets sont tous égaux à }\operatorname{card}\Lambda
\end{aligned}\tag*{(40)}
\]
% archive sheet number: 87
\scanpage{89}
On a donc une vraie structure polyédrale
\inserted{pour $P_{M,i}$} ssi $\Lambda=\Lambda_0$. Quand à $P_M$ tout entier,
c'est une vraie structure polyédrale ssi
\deleted{\unsure{on a}} de plus $\pi_M$ réduit à un élément,
i.e.
\[
\left\{\begin{aligned}
&P_M=(S_M,\unsure{A_M|F_M},F_M)\text{ est un }2\text{-polyèdre (i.e. à}\\
&\text{étoiles connexes) ssi on a}\\
&\Lambda\simeq\mathbb Z\text{ ou }\Lambda\simeq\mathbb Z/n\mathbb Z,\ \text{avec }n\in\{\deleted{\unsure{2}},3,4,6\}.\\
&\text{(NB le cas }n=2\text{ est exclu d'emblée ici !)}.
\end{aligned}\right.\tag*{(41)}
\]

En tous cas, on considère sur $S_M$ l'ens.
des structures triangulées $(F_{M,i})_{i\in I}$, \deleted{de}
\unsure{et|ce} sera cela pour nous $P_M$ (sans hypothèse
sur $\Lambda$ autre que $2\cdot1_\Lambda\ne0$). On aura
pour cette structure une suite exacte
\diagram{\(\text{\unsure{(42)|(22)}}\quad\begin{tikzcd}[ampersand replacement=\&,column sep=1.1em]1\arrow[r]\&\underbrace{\operatorname{Aut}_0(P_M)}_{\substack{\text{autom. de }P_M\\\text{fixant chaque}\\F_{M,i}}}\arrow[r]\&\operatorname{Aut}(P_M)\arrow[r,"\substack{\text{pas néc.}\\\text{surjectif}}"']\&\circlearrowright_{\pi_M}\end{tikzcd}\)}
et un hom de suites exactes
\diagram{\(\text{(43)}\quad\left\{\vcenter{\hbox{\begin{tikzcd}[ampersand replacement=\&,column sep=1.1em,row sep=3em]
1\arrow[r]\&{\unsure{\mathrm{Sl}|\mathrm{Gl}}}'_\Lambda(M)^{\pm}\arrow[r]\arrow[d]\&\mathrm{Gl}_\Lambda(M)'\arrow[r]\arrow[d]\&\deleted{\Lambda}Z'_\Lambda\arrow[r]\arrow[d,hook,"\substack{\text{op. sur }\pi_M\text{ qui}\\\text{est un tors-}\\\text{sous }Z'_\Lambda}"]\&1\\
1\arrow[r]\&\operatorname{Aut}_0(P_M)\arrow[r]\&\operatorname{Aut}(P_M)\arrow[r]\&\circlearrowright_{\pi_M}\&
\end{tikzcd}}}\right.\)}
où $Z'_\Lambda=\Lambda^*/\{\pm1_\Lambda\}$ déjà envisagé dans (13) (p.~36).
\deleted{\unsure{et} $Z'_\Lambda=$}
% archive sheet number: 88
\scanpage{90}\mspage{44}
\[
\begin{aligned}
&\text{Le groupe }\mathrm{Gl}'_\Lambda(M)^{\pm}\text{ est simplement transitif}\\
&\text{sur l'ens des repères }(s,t,u)\text{ de }P_{M,i}\ (i\in\pi\\
&\text{fixé), et }\mathrm{Sl}'_\Lambda(M)\text{ est simplement transitif}\\
&\text{sur l'ens des repères \uline{directs} de }P_{M,e}\\
&(e\in\overrightarrow\pi\text{ fixé})\ --\ \text{enfin, }\mathrm{Gl}'_\Lambda(M)^{\unsure{*|\prime|\varepsilon}}\text{ est}\\
&\text{simplement transitif sur l'ens des}\\
&\text{repères de }P_M.
\end{aligned}\tag*{(44)}
\]

\uline{Théorème} Supposons que $\Lambda=\Lambda_0$ i.e. que les
$P_{M,i}$ \inserted{$(i\in\pi_M)$} soient des vrais polyèdres. Alors
\[
\begin{aligned}
&\text{a) Pour } e\in\overrightarrow\pi,\qquad &&\mathrm{Sl}'_\Lambda(M)\xrightarrow{\sim}\operatorname{Aut}^+(P_{M,e})\\
&\text{b) Pour } i\in\pi,&&\mathrm{Gl}'_\Lambda(M)^{\pm}\xrightarrow{\sim}\operatorname{Aut}(P_{M,i})\\
&\text{c) \deleted{Pour} \unsure{On a}}&&\mathrm{Gl}'_\Lambda(M)\xrightarrow{\sim}\operatorname{Aut}(P_M)
\end{aligned}
\]
% slanted
\margin{\uline{NB} La connaissance de $F_M\subset\unsure{\mathfrak P|\mathbb P|\mathcal P}_3(S_M)$ permet déjà de reconstituer $A_M$ (33) et par la relation d'incidence naturelle, $(F_M,A_M)$ est un graphe strict, ayant $F_M$ comme ens. des sommets, $A_M$ comme ens des arêtes (les ``sommets'' sont d'ordre 3) -- ceci posé, la décomposition de $F_M$ en composantes $F_{M,i}$ est la suivante : c'est la décomposition \unsure{en composantes connexes|en composantes con-} du graphe précédent.}

où $P_M$ est défini, on le rappelle, comme la
structure formée par $S_M$ et l'ens des
$F_{M,i}$ de $\mathfrak P_3(S_M)$, i.e. l'ens des
structures de 2-polyèdre combinatoire régulier
sur l'ens des sommets $S_M$, associées aux
diverses structures $\frac12$ symplectiques sur $M$.

\uline{Scholie} On voit donc que l'on a des
équivalences \inserted{canoniques} de groupoïdes connexes ;
pour $\Lambda\simeq\mathbb Z$ ou $\Lambda\simeq\mathbb Z/n\mathbb Z$ fixé
(soit $\Lambda_n$)
% archive sheet number: 89
\scanpage{91}
\diagram{\(\text{(45)}\quad\begin{tikzcd}[ampersand replacement=\&,column sep=1.5em]
\parbox[t]{.32\linewidth}{Modules symplectiques libres de rang 2 sur $\Lambda$ ``modulo'' symétrie $\pm1$ (groupoïde quotient\ldots)}\arrow[r,"\approx"]\&
\parbox[t]{.38\linewidth}{groupoïde formé des 2-polyèdres \deleted{\unsure{triang.|triangulés}}\inserted{triangulés} orientés réguliers, isom. à un type \deleted{\unsure{tel}} bien déterminé (dépendant de $n$)}\\
\deleted{\unsure{\{}}(M,\inserted{e})\bmod-\mathrm{id}_M\arrow[r,mapsto]\&(S_M,F_{M,e}\unsure{\subset|\in}F_{M,\{e,-e\}})
\end{tikzcd}\)}

\diagram{\(\text{(46)}\quad\left\{\vcenter{\hbox{\begin{tikzcd}[ampersand replacement=\&,column sep=1.5em]
\parbox[t]{.34\linewidth}{Modules $\frac12$ symplectiques libres de rang 2 sur $\Lambda$ ``modulo'' symétrie $\pm1$}\arrow[r,"\approx"]\&
\parbox[t]{.36\linewidth}{\deleted{\unsure{conn.|connexes}} groupoïde formé des \deleted{certains} 2-polyèdres triangulés réguliers (orientables mais non orientés), de type déterminé (dépendant de $n$)}\\
(M,\deleted{i)\,\unsure{\bmod\pm}}\mathrm{id}_M\arrow[r,mapsto]\&(S_M,F_{M,i})
\end{tikzcd}}}\right.\)}

\diagram{\(\text{(47)}\quad\left\{\vcenter{\hbox{\begin{tikzcd}[ampersand replacement=\&,column sep=1.5em]
\parbox[t]{.34\linewidth}{Modules libres de rang 2 sur $\Lambda$, ``modulo'' symétrie $\pm1$}\arrow[r,"\approx"]\&
\parbox[t]{.36\linewidth}{groupoïde formé des \deleted{\unsure{struc.}}\inserted{structures} multi-2-polyédrales triangulées (paquet de structures 2-polyédr\unsure{ales} \unsure{\&} triangulées), de type d'iso. bien déterminé (dépendant de $n$)}\\
M\bmod-\mathrm{id}_M\arrow[r,mapsto]\&\bigl(S_M,(F_{M,i})_{i\in\pi_M}\bigr)
\end{tikzcd}}}\right.\)}

Il s'agirait alors

a) de caractériser axiomatiquement les structures ainsi obtenues

b) d'expliciter comment à une telle structure, on associe une $\mathbb F_2$-gerbe
(celle des modules $M$ dont la structure provient
% archive sheet number: 90
\scanpage{92}\mspage{45}
c) d'expliciter la construction des $M$
eux-mêmes \unsure{dont|\deleted{dont}} provient la structure
combinatoire envisagée

d) d'expliciter les autres structures qu'on
\unsure{trouve|tire|trouvent} \unsure{à} partir de la donnée combinatoire, telles : la structure de \deleted{$\Lambda^*$}
torseur sous $Z'_\Lambda=\Lambda^*/\pm1$ sur l'ens $\pi$,
le $\Lambda^*$-torseur $\overrightarrow\pi$ \inserted{qui relève $\pi$},
\deleted{l'application} (on
aura $\overrightarrow\pi=\coprod_{i\in\pi}\operatorname{Or}(P_{M,i})$ -- mais il faut
définir l'action de $\Lambda^*$ sur ce $\overrightarrow\pi$)
l'application $\overrightarrow\tau_F:\overrightarrow F_M\to\overrightarrow\pi$, l'opération
de $Z'_\Lambda$ sur $S_M$ \ldots

e) Expliciter comment l'ens des \deleted{\unsure{P}}structures
polyédrales $P_{M,i}$ se déduit de l'une
d'entre elles

f) Examiner, \inserted{pour $\Lambda$ quelconque,} les structures combinatoires
supplémentaires sur les $P_{M,i}$ ou $P_M$,
impliquant $\Lambda$ (p.~ex. l'opération de
$Z'_\Lambda$ sur $S_M$) qui permettent d'énoncer un théorème analogue
au précédent, et de réconstituer
$M$ à symétrie $-1$ près \ldots
% archive sheet number: 91
\scanpage{93}\mspage{46}
$\overline{\underline{\mathrm{VI}}}$ Je commencerai par la question e)
de la page précédente, qui avait déjà été
laissée en suspens au \S\,$\unsure{\mathrm{IV}|\mathrm{III}|\mathrm{VI}}$.

\deleted{\uline{Proposition}} Soit $\Lambda$ un anneau quelconque, \inserted{$M$ $\Lambda$-mod libre de rg 2}
considérons l'action de $\mathrm{Sl}_\Lambda(M)$ sur $M^*=\widetilde S_M$
de $\mathrm{Sl}_\Lambda(M)'$ sur $M^*/\pm1=S_M$. On cherche
les commutants de ces actions -- \deleted{dans les groupes $\unsure{\circlearrowright|\mathfrak S}_M$} \deleted{avec}
\inserted{dans les groupes des permutations de $\widetilde S_M$ resp. $\circlearrowright_M$}
présomption que c'est l'image de $\Lambda^*$
resp. de $Z'_\Lambda=\Lambda^*/\pm1$. Comme $M^*$ et a
fortiori \inserted{son quotient} $S_M$ \deleted{est} sont des espaces homogènes
sous \inserted{le Groupe d'\unsure{opération|opérateurs} $G$} $\mathrm{Sl}_\Lambda(M)$ resp. $\mathrm{Sl}_\Lambda(M)'$, le commutant
\inserted{dans $\circlearrowright_S$ de $G$} est isomorphe à $N_G(H)/H$, où $H$ est
le stabilisateur d'un point $s$ de $S$, de
façon précise on aura un \deleted{homomorphisme}
\deleted{\unsure{su|iso}}\inserted{isom}
\[
\deleted{\unsure{C}}Z\xrightarrow{\sim}N_G(H)/H
\]
par la relation entre $\deleted{\unsure{C}}Z$ et $N_G(H)$
définie par la formule
\[
u(s)=g\cdot s\qquad u\in Z,\ g\in N_G(H)\tag*{(48)}
\]

Calculons donc le normalisateur de
$H$, dans les deux cas, en prenant une base
$(x,y)$ de $M$, et en prenant $s=x$ ou $s=\dot x$.
% archive sheet number: 92
\scanpage{94}
A)
\[
\left\{\begin{aligned}
G&=\mathrm{Sl}(2,\Lambda)=\left\{\begin{pmatrix}a&b\\c&d\end{pmatrix}\in M_2(\Lambda)\mid ad-bc=1\right\}\\
H&=\left\{\begin{pmatrix}1&\deleted{\unsure{\lambda}}t\\0&1\end{pmatrix}\mid\deleted{\lambda}t\in\Lambda\right\}
\end{aligned}\right.\tag*{(49)}
\]

Soit $g$
\[
\begin{aligned}
g&=\begin{pmatrix}a&b\\c&d\end{pmatrix}\in G,\quad\text{d'où}\\
g^{-1}&=\begin{pmatrix}d&-b\\-c&a\end{pmatrix}\qquad\text{on a}
\end{aligned}
\]
\[
\operatorname{int}(g)\begin{pmatrix}1&\lambda\\0&1\end{pmatrix}=
\underbrace{\begin{pmatrix}a&b\\c&d\end{pmatrix}\begin{pmatrix}1&\deleted{\lambda}t\\0&1\end{pmatrix}}_{\begin{pmatrix}a&a\deleted{\lambda}t+b\\c&c\deleted{\lambda}t+d\end{pmatrix}}
\begin{pmatrix}d&-b\\-c&a\end{pmatrix}=
\begin{pmatrix}1\deleted{\unsure{+}}-ac\deleted{\lambda}t&a^2\deleted{\lambda}t\\-c^2\deleted{\lambda}t&1+ac\deleted{\lambda}t\end{pmatrix}\tag*{(*)}
\]
donc
\[
\operatorname{int}(g)\begin{pmatrix}1&\deleted{\lambda}t\\0&1\end{pmatrix}\in H
\Longleftrightarrow\left\{\begin{aligned}
-c^2\deleted{\lambda}t&=0\quad\text{d'où }(1-ac\lambda)(1+ac\lambda)=1\\
\deleted{\unsure{ac\lambda=0}}\\[-3pt]
1-ac\deleted{\lambda}t&=1
\end{aligned}\right.
\]
Donc
\[
N_G(H)=\left\{g=\begin{pmatrix}a&b\\c&d\end{pmatrix}\in\mathrm{Sl}(2,\Lambda)\ \middle|\ 
\begin{aligned}-c^2\deleted{\lambda}t&=0\\
\deleted{\unsure{1-}}ac\deleted{\lambda}t&=0\end{aligned}\right\}
\deleted{\unsure{\forall\lambda\in\Lambda\}}}
\]
Or \deleted{faisant $\lambda=1$, cela signifie $c^2=0$, \inserted{\unsure{et $ac$}} donc}
comme $ad-bc=1$ i.e. $ad=1+bc$, on voit
que \inserted{$c^2=0\Longrightarrow$} $ad$ est inversible \deleted{($(ad)^{-1}=1-bc$ avec} avec
\[
(ad)^{\unsure{-1|+1}}=1-bc\qquad a^{-1}=(1-bc)d
\]
et comme $ac=0$, on doit avoir $c=0$

Donc
\[
\begin{aligned}
N_G(H)&=\left\{g=\begin{pmatrix}\deleted{\unsure{a}}\lambda&b\\0&\lambda^{-1}\end{pmatrix}\mid\lambda\in\Lambda^*,\ b\in\Lambda\right\}\\
&\shortparallel\\[-3pt]
&\left\{g\in\mathrm{Sl}(2,\Lambda)\mid\operatorname{int}(g)\begin{pmatrix}1&1\\0&1\end{pmatrix}\in H\right\}
\end{aligned}\tag*{(50)}
\]
Donc une \deleted{\unsure{autom}} permutation de $M^*$ qui
commute à $\mathrm{Sl}(2,\Lambda)$ transforme $x$ en $\lambda x$, \inserted{$(\lambda\in\Lambda^*)$}
donc c'est l'homothétie $\lambda\,\mathrm{id}_M$, qui a le
$=$ effet sur $x$ et commute à $\mathrm{Sl}(2,\Lambda)$.
% archive sheet number: 93
\scanpage{95}\mspage{47}
B)
% slanted
\margin{On suppose maintenant $2\cdot1_\Lambda\ne0$}
\[
\left\{\begin{aligned}
G'&=\mathrm{Sl}(2,\Lambda)'=\deleted{\unsure{S}}G/\pm1\\
H'&=\text{image dans }G'\text{ du groupe }H\text{ de (49)}\\
&\hspace{3em}\text{par }p:G\to G'\quad\text{(can)}
\end{aligned}\right.\tag*{(51)}
\]
Soit $g'$ un élément de $G'$, provenant de
\[
g=\begin{pmatrix}a&b\\c&d\end{pmatrix},\qquad ad-bc=1
\]
et considérons \deleted{\unsure{pour $t$|pour le}} \deleted{les}
\[
\operatorname{int}(g')\cdot\underbrace{p\begin{pmatrix}1&t\\0&1\end{pmatrix}}_{h'=p(h)}
=p\ \operatorname{int}(g)\begin{pmatrix}1&t\\0&1\end{pmatrix}
=p\begin{pmatrix}1-act&a^2t\\-c^2t&1+act\end{pmatrix}
\]
% slanted
\margin{\uline{NB} l'image inv. de $H'$ dans $G$ est formé des $\left(\begin{smallmatrix}\pm1&\bullet\\0&\bullet\end{smallmatrix}\right)$}
on a
\[
\operatorname{int}(g')(h')\in H'\Longleftrightarrow
\left|\begin{aligned}-c^2t&=0\\1-act&\in\pm1\end{aligned}\right.\tag*{(52)}
\]
Donc
\[
g'\in N_{G'}(H')\Longleftrightarrow
\left\{\begin{aligned}&c^2=0\\&\forall t\in\Lambda,\ \text{on a}\quad act=0\text{ ou }act=2\end{aligned}\right.
\]

Considérons le cas $t=1$, on aura donc
\[
ac=0\quad\text{ou}\quad ac=2
\]
dans le premier cas, on a donc (puisque
$c^2=0\Longrightarrow a\in\Lambda^*$) $c=0$, ce qui exclut d'ailleurs
$ac=2$ (puisque $2_\Lambda\ne0$ par hypothèse). Dans
le deuxième cas, on doit avoir, en élevant
au carré
\[
4_\Lambda=0
\]
et de plus, la \inserted{$2^{\text e}$} relation (52) \deleted{\unsure{implique}} \inserted{s'écrit}
\[
1-2t\in\pm1\qquad\forall t\in\Lambda\qquad\text{i.e.}
\]
\[
\forall t\in\Lambda,\ \text{on a}\quad 2t=0\quad\text{ou}\quad 2t=2\tag*{(53)}
\]
ce qui implique que l'idéal $2\Lambda$ est réduit à deux
% archive sheet number: 94
\scanpage{96}
éléments $2_\Lambda$ et $0_\Lambda=4_\Lambda$. Notons que comme
$a^{-1}$ inv. et $2\ne0$, on a $c=a^{-1}\cdot2\ne0$ d'où $a^{-1}2=2$,
donc
\[
c=2
\]

Donc on a deux cas

1) Si \deleted{$\Lambda$ \unsure{est un anneau|est une} $4_\Lambda=0$, $2_\Lambda\ne0$, et $2\cdot\Lambda$ \unsure{réduit|est}} l'idéal
$2\Lambda$ n'a que deux éléments $0$ et $2$ (i.e.
$\Lambda$ est une extension d'une $\mathbb F_2$-algèbre \inserted{$\Lambda_0$} par
un idéal engendré par $\unsure{2\cdot1|2_\Lambda|2}$ et ayant $2$ comme
seul élément non nul) alors
% slanted
\margin{Ce cas a lieu où $\Lambda$ est l'image inverse de l'extension $\mathbb Z/4\mathbb Z$ de $\mathbb F_2$ par $\mathbb F_2$, par un hom de $\mathbb F_2$-alg. $\Lambda_0\to\mathbb F_2$}
\[
N_{G'}(H')=\left\{p\begin{pmatrix}a&b\\c&d\end{pmatrix}\ \middle|\ 
\begin{gathered}\deleted{\unsure{a\ne0}},ad-bc=1\\c\in\{0,2\}\end{gathered}\right\}\tag*{(54)}
\]
admet $N^0_{G'}(H')\overset{\text{déf}}{=}\left\{p\left(\begin{smallmatrix}\lambda&b\\0&\lambda^{-1}\end{smallmatrix}\right)\mid\lambda\in\Lambda^*,b\in\Lambda\right\}$
comme sous-groupe d'indice $2$.
% slanted
\margin{en fait, si $\sigma=p\left(\begin{smallmatrix}1&0\\2&1\end{smallmatrix}\right)$ alors $N_{G'}(H')$ est produit \inserted{$1/2$ direct} de $N^0_{G'}(H')$ par $\{1,\sigma\}$. [Le $\sigma$ ou antipodisme a un sens intrinsèque, indépend$^{\text{mt}}$ du choix d'une base}

2) Dans tout autre cas, \deleted{(p. ex. si $4_\Lambda\ne0$),}
\[
N_{G'}(H')=N^0_{G'}(H')\tag*{(55)}
\]

Dans ce cas \inserted{$2^\circ$}, on voit donc que le commutant
de $\mathrm{Sl}_\Lambda(M)'$ dans $\unsure{S|\widetilde S|\mathfrak S}_M=M^*/\pm1$ est $Z'_\Lambda$.
Dans le cas $1^\circ$) par contre, $Z'_\Lambda$ est d'indice
$2$ dans ce commutant, qui est isom. au
produit direct \inserted{(pas seulement $1/2$ direct!)} de $Z'_\Lambda$ par le sous-groupe
$\{1,\sigma\}$. Le cas typique est celui de
$\Lambda=\mathbb Z/4\mathbb Z$, \deleted{$\sigma(\dot x)=(x+2y)^{\displaystyle\cdot}={}$} alors $Z'_\Lambda=\{1\}$,
mais le groupe des autom de $S_M$ (ens des
sommets de l'octaèdre $P_M$) qui commutent à
$\operatorname{Aut}^{+}(P_M)$, n'est pas réduit à $\{1\}$, il est engendré
par l'antipodisme de l'octaèdre.
% archive sheet number: 95
